\chapter{Reproducibilidad y comandos de referencia}
\label{app:reproducibilidad}

\section{Ambiente reportado}
Los reportes de cierre indican ejecución principal en plataforma Windows 11 con 12 núcleos lógicos, Python 3.14.x, NumPy 2.4.x, Pandas 3.0.x y SciPy 1.17.x. En la práctica, los números exactos de versión deben congelarse nuevamente antes de una entrega final, idealmente mediante \code{requirements.txt}, \code{uv.lock} o un contenedor reproducible.

\section{Comandos base}
La corrida B2/B3/B4 se documentó mediante comandos de referencia equivalentes a:
\begin{verbatim}
$env:PYTHONPATH='.../Thesis-Copilot-Toolkit'
$env:INCLUDE_MNE='1'
$env:MAX_TIME_SAMPLES='220'
$env:B2_DTW_MAX_POINTS='80'
$env:B2_TOP_K='3'
$env:B2_RANDOM_SEED='42'
python experiments/run_b2_batched.py
python experiments/consolidate_b2_publication.py
python experiments/finalize_b3_b4.py
\end{verbatim}

\section{Criterios GO/NO-GO}
El cierre documental distingue resultados confirmatorios de iteraciones exploratorias. La regla empleada es que una iteración \textit{NO-GO} puede discutirse como evidencia de sensibilidad, fallo o limitación, pero no debe usarse como soporte principal de una afirmación positiva. El caso INS-13/BGSRP se reporta explícitamente como validación proxy o parcial, no como equivalencia estricta 1:1 con MATLAB/GSPBox.

\section{Checklist mínimo antes de entrega institucional}
\begin{enumerate}
\item Compilar el documento desde la carpeta \code{hermes/tesis\_magister\_gsp\_eeg\_hermes}.
\item Verificar que no existan citas ``?'' ni referencias cruzadas sin resolver.
\item Actualizar examinadores, fecha, agradecimientos y cualquier campo institucional pendiente.
\item Confirmar que las figuras copiadas en \code{figures/} corresponden a los artefactos finales.
\item Revisar que las tablas de apéndice deriven de resultados compatibles en normalización.
\item Si se entrega a revisión externa, congelar commit, ambiente y archivo PDF final.
\end{enumerate}

\section{Benchmark confirmatorio contra MNE-Python}
La comparación final TRSS vs. MNE se ejecutó y analizó con un ambiente reproducible basado en \code{uv}, ya que el Python del sistema no contenía todas las dependencias científicas. El comando de análisis estadístico robusto fue:
\begin{verbatim}
uv run --python 3.12 --with pandas --with numpy --with scipy --with tabulate \
  python experiments/analyze_trss_mne_robust_stats.py
\end{verbatim}

El benchmark extensivo se encuentra en:
\begin{verbatim}
experiments/trss_vs_mne_bads_extensive.py
results/trss_vs_mne_bads_extensive/
\end{verbatim}

La versión de MNE inspeccionada durante esta etapa fue \code{1.12.1}. La implementación de referencia usa interpolación por splines esféricas para EEG y se consultó tanto desde la documentación pública de MNE-Python como desde el código fuente instalado localmente.

\section{Compilación LaTeX en ambiente sin privilegios}
La compilación se realizó con TinyTeX instalado en modo usuario, porque la instalación por \code{sudo apt} no estaba disponible en el entorno WSL. El comando de validación recomendado es:
\begin{verbatim}
latexmk -pdf -interaction=nonstopmode -halt-on-error tesis_magister_hermes.tex
\end{verbatim}

